\lecture{9.9}{Rubin's Missingness Hierarchy}
\label{lec:9-9}

\subsection{Missing Completely at Random}

\begin{definition}[Missing completely at random]
Missingness is missing completely at random (MCAR) if the probability of
missingness does not depend on any data values, observed or unobserved.
\[
  R\ci (X_{\obs},X_{\mis},Z),
  \qquad
  p(R\given X,Z)=p(R).
\]
\end{definition}

MCAR is like randomly removing entries from a data set. Complete cases are a
random subsample with respect to the variables in the full data. This is a strong
condition and often implausible in medical and social data.

\subsection{Missing at Random}

\begin{definition}[Missing at random]
Missingness is missing at random (MAR) if the probability of missingness may
depend on observed information but not on missing values after conditioning on
the observed information:
\[
  R\ci X_{\mis}\given Z,X_{\obs},
  \qquad
  p(R\given X,Z)=p(R\given Z,X_{\obs}).
\]
\end{definition}

For example, outcome missingness may depend on age and sex if those variables
are observed for everyone. Under MAR, methods such as inverse probability
weighting and multiple imputation can be justified if the models are correctly
specified.

\subsection{Missing Not at Random}

\begin{definition}[Missing not at random]
Missingness is missing not at random (MNAR) if the probability of missingness
depends on missing values themselves, even after conditioning on observed
information.
\end{definition}

For example, income may be missing because high-income respondents choose not
to report it, or substance-use history may be missing because people with
substance-use history refuse to answer. MNAR is not one model; it is a broad
category containing many different mechanisms.

\subsection{Statistical Definitions Using $R$, $X_{\text{obs}}$, $X_{\text{miss}}$, and $Z$}

Let $Z$ denote variables observed for all units. Let $X$ denote variables subject
to missingness, decomposed into observed and missing parts:
\[
  X=(X_{\obs},X_{\mis}).
\]
Let $R$ denote the missingness pattern. The hierarchy can be summarized as:
\[
  \text{MCAR:}\qquad p(R\given X,Z)=p(R),
\]
\[
  \text{MAR:}\qquad p(R\given X,Z)=p(R\given X_{\obs},Z),
\]
\[
  \text{MNAR:}\qquad p(R\given X,Z)
  \text{ depends on }X_{\mis}\text{ for at least one pattern}.
\]

\begin{figure}[H]
\centering
\begin{tikzpicture}[
  box/.style={draw, rounded corners=2pt, minimum width=2.5cm, minimum height=.75cm},
  arr/.style={-{Latex[length=2.2mm]}, thick},
  node distance=.9cm
]
\node[box] (mcar) {MCAR\\ $R$ independent of data};
\node[box, below=of mcar] (mar) {MAR\\ $R$ depends on observed data};
\node[box, below=of mar] (mnar) {MNAR\\ $R$ depends on missing data};
\draw[arr] (mcar) -- (mar);
\draw[arr] (mar) -- (mnar);
\end{tikzpicture}
\caption{Rubin's hierarchy classifies missingness by what information predicts missingness.}
\label{fig:ch9-rubin-hierarchy}
\end{figure}

\subsection{Limitations of MCAR/MAR/MNAR}

The hierarchy is useful but incomplete. Its limitations include:
\begin{itemize}
  \item MCAR and MAR are statistical restrictions, not causal explanations.
  \item MNAR is too broad: many distinct mechanisms are all labeled MNAR.
  \item The hierarchy can be hard to apply when several variables are missing.
  \item It does not automatically tell us whether a target estimand is
  recoverable.
  \item It does not by itself specify an estimation strategy.
\end{itemize}

\begin{keyidea}
MCAR, MAR, and MNAR describe how missingness probabilities relate to observed
and unobserved values. They are useful labels, but causal graphs provide a more
explicit language for why values are missing and what can be recovered.
\end{keyidea}
